% Source-bound notation snippets; reference only. No macro definitions or global renaming.
% Requires bm only when a chosen snippet uses \bm. Existing manuscript definitions take priority.
% index_set: 从1到N的有限索引集
% $[N]$
% Source: dr_MDP.tex:L166-L166
% ones_vector: 全一向量
% $\bm{e}$
% Source: dr_MDP.tex:L166-L166
% basis_vector: 第j个坐标为1、其余为0的向量
% $\bm{e}_j$
% Source: dr_MDP.tex:L166-L166
% simplex: 概率单纯形
% $\Delta_N$
% Source: dr_MDP.tex:L166-L166
% probability_measures: 集合X上的全部Borel概率测度
% $P(\mathcal{X})$
% Source: dr_MDP.tex:L166-L166
% dirac: 在向量x处的狄拉克分布
% $\delta_{\bm{x}}$
% Source: dr_MDP.tex:L166-L166
% state_set: 有限状态集合
% $\mathcal{S}$
% Source: dr_MDP.tex:L220-L220
% state_count: 状态数
% $S$
% Source: dr_MDP.tex:L220-L220
% action_set: 有限动作集合
% $\mathcal{A}$
% Source: dr_MDP.tex:L220-L220
% action_count: 动作数
% $A$
% Source: dr_MDP.tex:L220-L220
% initial_distribution: 初始状态分布
% $\bm{p}_0$
% Source: dr_MDP.tex:L220-L220
% transition_vector: 给定状态动作后的下一状态分布
% $\bm{p}_{sa}$
% Source: dr_MDP.tex:L220-L220
% transition_block: 固定状态下全部动作的转移核块
% $\bm{p}_s$
% Source: dr_MDP.tex:L220-L220
% transition_kernel: 全部状态动作的转移核
% $\bm{p}$
% Source: dr_MDP.tex:L220-L220
% transition_reward: 状态s执行动作a转到s′的奖励
% $r_{sas^{\prime}}$
% Source: dr_MDP.tex:L220-L220
% reward_vector: 固定状态动作下按下一状态汇集的奖励
% $\bm{r}_{sa}$
% Source: dr_MDP.tex:L220-L220
% discount_factor: 折扣因子
% $\lambda$
% Source: dr_MDP.tex:L220-L220
% policy: 随机策略
% $\bm{\pi}$
% Source: dr_MDP.tex:L226-L226
% state_policy: 给定状态的动作分布
% $\bm{\pi}_s$
% Source: dr_MDP.tex:L226-L226
% policy_probability: 动作概率分量
% $\pi_{s a}$
% Source: dr_MDP.tex:L248-L248
% value_vector: 状态价值向量及其迭代、最优标记
% $\bm{v}$
% Source: dr_MDP.tex:L260-L264
% bellman_operator: 分布鲁棒贝尔曼算子
% $\mathfrak{T}$
% Source: dr_MDP.tex:L260-L264
% kernel_law: 全转移核的未知概率分布
% $\mu$
% Source: dr_MDP.tex:L234-L234
% state_kernel_law: 状态核块的边际分布
% $\mu_s$
% Source: dr_MDP.tex:L236-L236
% law_ambiguity: 状态矩形的转移核分布歧义集
% $\mathbb{M}^q$
% Source: dr_MDP.tex:L234-L234
% state_law_ambiguity: 以Wasserstein球限制状态核分布的集合
% $\mathbb{M}_{s}^q$
% Source: dr_MDP.tex:L238-L238
% reference_law: Wasserstein球中心的参考/经验分布
% $\nu_s$
% Source: dr_MDP.tex:L240-L246
% kernel_sample_count: 经验分布中的转移核样本数
% $N$
% Source: dr_MDP.tex:L246-L246
% sample_kernel_block: 第i个状态核样本块
% $\hat{\bm{p}}^i_s$
% Source: dr_MDP.tex:L246-L246
% perturbed_kernel_vector: 样本i对应的可选/扰动转移概率
% $\bm{p}^i_{sa}$
% Source: dr_MDP.tex:L274-L279
% sample_kernel_vector: 第i个样本在状态动作处的转移向量
% $\hat{\bm{p}}^{i}_{sa}$
% Source: dr_MDP.tex:L274-L279
% wasserstein_radius: Wasserstein歧义球的半径
% $\theta$
% Source: dr_MDP.tex:L238-L238
% wasserstein_order: Wasserstein阶数及本文地面范数阶数
% $q$
% Source: dr_MDP.tex:L238-L244
% wasserstein_distance: q阶Wasserstein距离
% $W_q$
% Source: dr_MDP.tex:L242-L244
% ground_distance: 用于构造Wasserstein距离的地面距离
% $d_q$
% Source: dr_MDP.tex:L244-L244
% coupling: 两个状态核分布之间的耦合分布
% $\kappa$
% Source: dr_MDP.tex:L242-L244
% coupling_set: 给定两边际分布的全部耦合
% $\Gamma$
% Source: dr_MDP.tex:L244-L244
% expected_kernel_set: Wasserstein分布歧义集导出的期望核集合
% $\mathbb{B}_s^q$
% Source: dr_MDP.tex:L250-L256
% bellman_threshold: Bellman重构的目标上界与外层二分量
% $\gamma$
% Source: dr_MDP.tex:L274-L279
% dual_multiplier: 内层约束的非负对偶乘子
% $\alpha$
% Source: dr_MDP.tex:L317-L324
% reward_plus_value: 奖励向量加折扣后的当前价值向量
% $\bm{b}_{sa}$
% Source: dr_MDP.tex:L290-L290
% projection_subproblem: 固定动作的最小扰动预算子问题
% $\mathfrak{P}$
% Source: dr_MDP.tex:L283-L290
% inner_subproblem: 单样本单动作的内层单纯形子问题
% $\mathfrak{D}_q$
% Source: dr_MDP.tex:L322-L324
% outer_tolerance: 外层γ二分的停止精度
% $\epsilon_1$
% Source: dr_MDP.tex:L374-L395
% inner_tolerance: 内层α二分的停止精度
% $\epsilon_2$
% Source: dr_MDP.tex:L374-L395
% budget_accumulator: 按动作累加的已计算预算/子问题值
% $r_a$
% Source: dr_MDP.tex:L374-L395
% dual_slope: 内层对偶目标的斜率
% $m$
% Source: dr_MDP.tex:L385-L386
% policy_set: 全部平稳随机策略集合
% $\Pi$
% Source: ho_petrik_wiesemann_phi_original.tex:L132-L132
% kernel_ambiguity: 鲁棒MDP的转移核歧义集
% $\mathcal{U}$
% Source: dr_MDP.tex:L228-L228
% phi_state_set: φ散度鲁棒模型中的逐状态核歧义集
% $\mathcal{P}_s$
% Source: ho_petrik_wiesemann_phi_original.tex:L159-L159
% phi_budget: φ散度模型的不确定性预算
% $\kappa$
% Source: ho_petrik_wiesemann_phi_original.tex:L161-L161
% phi_generator: φ散度的凸生成函数
% $\phi$
% Source: ho_petrik_wiesemann_phi_original.tex:L165-L165
